\scanpage{88}
% archive sheet number: 90

\[
X(I)/V(I)\simeq X(J)\qquad\text{où }J=J(I)\text{ comme dessus}
\]

Ceci dit soit \deleted{$I\hookrightarrow X$ on \unsure{a donc|se donne} donc dans $X$ une opération $V=V(I)$}\inserted{une opération de $V$ sur $X$} $V\hookrightarrow G_X$, d'où un $I(V)$ $V$-torseur des bi-faces de $\mathrm{Oc}(V)$, et on a \deleted{$X\simeq$}

\[
X\simeq X(I(V))\qquad\text{iso. canonique}
\]

(et non $X\simeq X(I)$). Mais pour \deleted{un} deux torseurs $T,T'$ sous un $=$ $V$, on aura

\[
X(T\wedge T')=X(T)\wedge_V T'
\]

\margin{à expliciter}

Soit donc
\[
\begin{aligned}
\underline{t}=\delta(I,\unsure{X|\lambda})&\text{ le }V\text{-torseur tel que }I(V)\simeq I\wedge_V\tau\\
&\text{i.e. }\tau\simeq I(V)\wedge I^{-1}
\end{aligned}
\]

on aura donc

\[
\boxed{X=X(I(V))=X(I)\wedge_V\underline t}
\]

\deleted{Donc} \deleted{et l'on a} \inserted{donc}

\[
\boxed{X/V\simeq X(I)/V\simeq \deleted{X}Z(J)}
\]

En résumé :

\uline{\unsure{bis|Loi|Cor.} : Prop} \deleted{La donnée, d'une part $V=\mathbb D_2$} \inserted{Pour $I\in\operatorname{Ob}\mathbb E\mathrm{ns}_4$} donné, d'où $J$ et $V=V(I)$, la donnée de un $X$ et d'un plongement $V\subset \deleted{\unsure{V_X|G}}G_X$ équivaut à la donnée d'un torseur $\underline t$ sous $V$,

\[
\text{via}\qquad \tau\longmapsto X(I)\wedge_V\deleted{\tau}\underline t
\]

\[
\text{On a alors}\qquad X/V\simeq\deleted{X}Z(J)\qquad\text{(iso. canonique).}
\]

\deleted{\uline{Corollaire}} \inserted{Donc} la donnée de $(X,I\hookrightarrow X)$ équivaut à celle de a) un pts $\zeta$ \unsure{de|/ de} $Z(J)^*$, b) un $V$-torseur $\underline t$, \deleted{c) Un $V$-isom \unsure{un|tel} tel que $X=X(I)\wedge_V\underline t$, et} \deleted{un \unsure{espace}.} \deleted{$X_\zeta=$ fibre} de \unsure{.} $X_\zeta=X(I)_\zeta\wedge\underline t\simeq I$ i.e. $\underline t=I\wedge X(I)_\zeta^{-1}$
